\section{Conclusion}
We developed an exact SAT pipeline for minimum-span $L(h,k)$-labeling,
established the correctness of its encoding and span search, and evaluated
it against assignment ILP on a fixed multi-family cohort. In the repeated
30-second comparison, SAT and ILP establish optimality in all three runs
on \ProgressSATStable{} and \ProgressILPStable{} of 117 instances,
respectively. The remaining-bound analysis shows why feasible-solution
quality and optimality proofs must be evaluated separately.

Combinatorial postprocessing complements this comparison. Combining the
existing evidence with verified square-clique bounds and explicit
color-class ordering witnesses closes \ConstructionClosed{} instances
and leaves \ConstructionOpen{} open. Of the closed instances,
\ConstructionCliqueCertified{} have matching square-clique bounds and
labelings, allowing a solver-independent optimality check. These
certificates concern the evaluated instances and are not additional
timed solver successes or a new exact formula for an entire graph class.

The contribution is a correct general-graph formulation with an auditable
comparison of exact methods and explicit limits of their current evidence.
The study does not claim that order encoding is new or that SAT dominates
ILP. A controlled encoding ablation and independently checkable UNSAT
traces are possible future extensions. Resolving every remaining interval
is not required to report the established results and computational limits
of the present experiment.
